\documentclass[../../main.tex]{subfiles}% 注意这里的文件路径不能用 ./main.tex ,否则用latexmk编译子文件会报错
\graphicspath{{\subfix{./image/}}} % 指定图片目录,后续可以直接使用图片文件名
% 注意这里的文件路径不能用 ../../image/ ,否则用latexmk编译子文件会报错

% 例如:
% \begin{figure}[H]
% \centering
% \includegraphics[scale=0.3]{图.png}
% \caption{}
% \label{figure:图}
% \end{figure}
% 注意:上述\label{}一定要放在\caption{}之后,否则引用图片序号会只会显示??.

\begin{document}

\section{主理想整环上的有限生成扭模}

本节总假定$D$是主理想整环（p.i.d.）。

\begin{lemma}
设$M$是$D$上有限生成的扭模，定义$M(a)=\{x\in M\mid ax=0\},\forall a\in D$，则下列结论成立：
\begin{enumerate}[(1)]
\item $\forall a\in D$，$M(a)$是$M$的子模；
\item $M(0)=M$，若$a$可逆，则$M(a)=\{0\}$；
\item 若$a\mid b$，则$M(a)\subseteq M(b)$；
\item 设$a,b\in D$，则$M(a)\cap M(b)=M((a,b))$，其中$(a,b)$为$a$与$b$的最大公因子；
\item 若$(a,b)=1$，则$M(ab)=M(a)\oplus M(b)$。
\end{enumerate}
\end{lemma}
\begin{proof}
\begin{enumerate}[(1)]
\item 任取$x,y\in M(a),c\in D$，有
\begin{align*}
a(x-y)=ax-ay=0,\quad a(cx)=c(ax)=0,
\end{align*}
故$M(a)$是$M$的子模。

\item 显然$M(0)=M$。若$a$可逆，对任意$x\in M(a)$，有$x=1\cdot x=a^{-1}(ax)=0$，故$M(a)=\{0\}$。

\item 若$a\mid b$，则$\exists c\in D$使得$b=ac$。任取$x\in M(a)$，有$bx=cax=0$，故$x\in M(b)$，即$M(a)\subseteq M(b)$。

\item 因$(a,b)\mid a,(a,b)\mid b$，由(3)得$M(a)\cap M(b)\supseteq M((a,b))$。
又$D$为主理想整环，故$\exists u,v\in D$使得
\begin{align*}
(a,b)=ua+vb.
\end{align*}
任取$x\in M(a)\cap M(b)$，有$ax=bx=0$，于是
\begin{align*}
(a,b)x=uax+vbx=0,
\end{align*}
故$x\in M((a,b))$。因此$M(a)\cap M(b)=M((a,b))$。

\item 由$(a,b)=1$，结合(4)得$M(a)\cap M(b)=M(1)=\{0\}$，即$M(a)$与$M(b)$无关。
又$a\mid ab,b\mid ab$，故$M(ab)\supseteq M(a)\oplus M(b)$。
$\exists u,v\in D$使得$au+bv=1$。任取$x\in M(ab)$，有
\begin{align*}
x=aux+bvx,
\end{align*}
且$b(aux)=0,a(bvx)=0$，故$x\in M(a)\oplus M(b)$。因此$M(ab)=M(a)\oplus M(b)$。
\end{enumerate}

\end{proof}

\begin{theorem}
设$M$是$D$上的有限生成扭模，$a\in D,a\neq0$，其素因式分解为
\begin{align*}
a=up_1^{n_1}p_2^{n_2}\cdots p_r^{n_r},
\end{align*}
其中$u$为$D$中可逆元，$p_1,p_2,\cdots,p_r$为互不相伴的元素，则
\begin{align*}
M(a)=M(p_1^{n_1})\oplus M(p_2^{n_2})\oplus\cdots\oplus M(p_r^{n_r}).
\end{align*}
\end{theorem}
\begin{proof}

对素因子个数$r$用数学归纳法证明。
当$r=1$时，$a=up_1^{n_1}$。因$a\mid p_1^{n_1},p_1^{n_1}\mid a$，由引理的(3)得$M(a)=M(p_1^{n_1})$。

设定理对$r-1$成立，下证对$r$成立。令
\begin{align*}
a_1=up_1^{n_1}p_2^{n_2}\cdots p_{r-1}^{n_{r-1}},\quad b=p_r^{n_r},
\end{align*}
则$a_1b=a,(a_1,b)=1$。由引理的(5)得$M(a)=M(a_1)\oplus M(b)$。
由归纳假设
\begin{align*}
M(a_1)=\bigoplus\limits_{i=1}^{r-1}M(p_i^{n_i}),\quad M(b)=M(p_r^{n_r}),
\end{align*}
因此
\begin{align*}
M(a)=\bigoplus\limits_{i=1}^{r}M(p_i^{n_i}).
\end{align*}

\end{proof}

\begin{definition}
设$M$是$D$-模，$p$是$D$的元素，称$M$的子集
\begin{align*}
M_p=\bigcup\limits_{i=1}^{\infty}M(p^i)
\end{align*}
为$M$的\textbf{$p$-分支}。若$M_p=M$，则称$M$为\textbf{$p$-模}。
\end{definition}
\begin{remark}
由引理的(3)，有
\begin{align*}
M(p)\subseteq M(p^2)\subseteq\cdots\subseteq M(p^k)\subseteq\cdots,
\end{align*}
故$M_p$是$M$的子模。
\end{remark}

\begin{definition}
设$R$是幺环，$M$是$R$-模。$R$的子集
\begin{align*}
\mathrm{ann}M=\{a\in R\mid ax=0,\forall x\in M\}
\end{align*}
称为$M$的\textbf{零化子}。
\end{definition}
\begin{remark}
$\mathrm{ann}M$是$R$的理想，且$\mathrm{ann}M=\bigcap\limits_{x\in M}\mathrm{ann}x$。
\end{remark}

\begin{lemma}
设$M$是$D$上的扭模，则
\begin{enumerate}[(1)]
\item $\forall x\in M$，$\exists a_x\in D$使得$\mathrm{ann}x=\langle a_x\rangle$；且$\exists a_M\in D$，使得$\mathrm{ann}M=\langle a_M\rangle$，其中$a_M$是$\{a_x\mid x\in M\}$的最小公倍式；
\item 若$M\neq\{0\}$且$M=\langle x_1,x_2,\cdots,x_r\rangle$，则$a_M=[a_{x_1},a_{x_2},\cdots,a_{x_r}]$，即$a_M$是$a_{x_1},a_{x_2},\cdots,a_{x_r}$的最小公倍式，且$\mathrm{ann}M$是$D$的非平凡理想。
\end{enumerate}
\end{lemma}
\begin{proof}
\begin{enumerate}[(1)]
\item 因$D$为主理想整环，$\mathrm{ann}x$为$D$的理想，故$\exists a_x\in D$使得$\mathrm{ann}x=\langle a_x\rangle$。
由$\mathrm{ann}M\subseteq\mathrm{ann}x$得$a_x\mid a_M,\forall x\in M$。
若$c\in D$满足$a_x\mid c,\forall x\in M$，则$c\in\mathrm{ann}x,\forall x\in M$，故
\begin{align*}
c\in\bigcap\limits_{x\in M}\mathrm{ann}x=\mathrm{ann}M=\langle a_M\rangle,
\end{align*}
因此$a_M\mid c$，即$a_M$是$\{a_x\mid x\in M\}$的最小公倍式。

\item 由(1)，$a_M$是$a_{x_1},a_{x_2},\cdots,a_{x_r}$的公倍式。
若$c\in D$满足$a_{x_i}\mid c,\forall i$，对任意$x\in M$，$\exists b_i\in D$使得$x=\sum\limits_{i=1}^{r}b_ix_i$，则
\begin{align*}
cx=c\left(\sum\limits_{i=1}^{r}b_ix_i\right)=\sum\limits_{i=1}^{r}b_icx_i=0,
\end{align*}
故$c\in\langle a_M\rangle$，即$a_M\mid c$，因此$a_M=[a_{x_1},a_{x_2},\cdots,a_{x_r}]$。
又$M\neq\{0\}$，$a_M$不可逆，故$\langle a_M\rangle\neq D$；且$a_M\neq0$，因此$\mathrm{ann}M$是$D$的非平凡理想。
\end{enumerate}

\end{proof}

\begin{theorem}
设$M$是主理想整环$D$上有限生成的扭模，$\mathrm{ann}M=\langle a\rangle$，其中
\begin{align*}
a=up_1^{n_1}p_2^{n_2}\cdots p_r^{n_r},\quad n_i\geqslant1,\ 1\leqslant i\leqslant r,
\end{align*}
$u$为$D$中可逆元，$p_i$为互不相伴的元素，则
\begin{enumerate}[(1)]
\item 对$D$中元素$p$，有
\begin{align*}
M_p=
\begin{cases}
M(p_i^{n_i}),& p\text{与}p_i\text{相伴},\ 1\leqslant i\leqslant r,\\
\{0\},& p\text{与}p_i\text{不相伴},\ 1\leqslant i\leqslant r;
\end{cases}
\end{align*}
\item $M=M_{p_1}\oplus M_{p_2}\oplus\cdots\oplus M_{p_r}$。
\end{enumerate}
\end{theorem}
\begin{proof}

由$\mathrm{ann}M=\langle a\rangle$，$\forall x\in M$有$ax=0$，即$M(a)=M$。结合定理5.4.1得
\begin{align*}
M=M(a)=\bigoplus\limits_{i=1}^{r}M(p_i^{n_i}).
\end{align*}

任取$p\in D,k\in\mathbb{N}$。若$p$与任一$p_i$不相伴，则$(p^k,a)=1$，于是
\begin{align*}
M(p^k)=M(p^k)\cap M=M((p^k,a))=M(1)=\{0\},
\end{align*}
因此
\begin{align*}
M_p=\bigcup\limits_{k=1}^{\infty}M(p^k)=\{0\}.
\end{align*}

若$p$与$p_i$相伴，则对任意$k\in\mathbb{N}$，$(p^k,a)=p_i^{\min\{k,n_i\}}$，于是
\begin{align*}
M(p^k)=M(p_i^{\min\{k,n_i\}})\subseteq M(p_i^{n_i}),
\end{align*}
故
\begin{align*}
M_p=\bigcup\limits_{k=1}^{\infty}M(p_i^k)=M(p_i^{n_i}).
\end{align*}
综上结论成立。

\end{proof}

\begin{corollary}
设$N$为$M$的子模，则
\begin{align*}
N=\bigoplus\limits_{i=1}^{r}N_{p_i},\quad N_{p_i}=N\cap M_{p_i}.
\end{align*}
\end{corollary}
\begin{proof}

设$\mathrm{ann}N=\langle b\rangle$，由$N\subseteq M$得$\mathrm{ann}M\subseteq\mathrm{ann}N$，故$b\mid a$。于是
\begin{align*}
b=u_1p_1^{k_1}p_2^{k_2}\cdots p_r^{k_r},\quad k_i\leqslant n_i,\ 1\leqslant i\leqslant r.
\end{align*}
由定理5.4.2得
\begin{align*}
N=\bigoplus\limits_{i=1}^{r}N(p_i^{k_i})=\bigoplus\limits_{i=1}^{r}N_{p_i},\quad N(p_i^{k_i})\subseteq M(p_i^{n_i}),
\end{align*}
故$N_{p_i}\subseteq M_{p_i}\cap N$。又$M_{p_i}\cap N\subseteq N_{p_i}$，因此$N_{p_i}=N\cap M_{p_i}$。

\end{proof}

\begin{lemma}
设$D$是主理想整环，$M$是$D$上秩为$n$的自由模，$N$是$M$的子模，商模$M/N$是$p$-模，则存在$M$的基$e_1,e_2,\cdots,e_n$及非负整数组$k_1,k_2,\cdots,k_n$，使得$p^{k_1}e_1,p^{k_2}e_2,\cdots,p^{k_n}e_n$为$N$的一组基。
\end{lemma}
\begin{proof}

对$M$的秩$r(M)$用数学归纳法证明。

当$r(M)=1$时，设$e$为$M$的基。因$M/N$是有限生成模，$\mathrm{ann}(M/N)=\langle p^k\rangle$，故$p^ke\in N$。
设$\pi:M\to M/N$为自然同态，若$a\in D$满足$ae\in N$，则$\pi(ae)=a\pi(e)=0$，于是$p^k\mid a$，故$N=D(p^ke)$，即$p^ke$为$N$的基。

设秩为$n-1$时引理成立，下证秩为$n$时成立。
任取$M$的一组基$f_1,f_2,\cdots,f_n$，定义
\begin{align*}
I_i=\{a_i\in D\mid \text{若}x=\sum\limits_{j=1}^{n}a_jf_j\in N\},\quad i=1,2,\cdots,n,
\end{align*}
则$I_i$为$D$的理想。因$D$为主理想整环，$\exists c_i\in D$使得$I_i=\langle c_i\rangle$。
又$M/N$为$p$-模，对$f_i$，$\exists n_i\in\mathbb{N}$使得$p^{n_i}f_i\in N$，故$p^{n_i}\in\langle c_i\rangle$，即$\exists k_i\leqslant n_i$使得$c_i$与$p^{k_i}$相伴，即$I_i=\langle p^{k_i}\rangle$。
不妨设$k_1=\min\{k_i\mid 1\leqslant i\leqslant n\}$，则$\exists g\in N$使得
\begin{align*}
g=p^{k_1}f_1+a_2f_2+\cdots+a_nf_n.
\end{align*}
因$a_i\in I_i=\langle p^{k_i}\rangle$，故$a_i=p^{k_i}a_i',\ 2\leqslant i\leqslant n$。令
\begin{align*}
e_1=f_1+\sum\limits_{j=2}^{n}a_j'f_j,
\end{align*}
则$e_1,f_2,\cdots,f_n$是$M$的基，且$p^{k_1}e_1=g\in N$。
令$M_1=\bigoplus\limits_{j=2}^{n}Df_j$，$N_1=N\cap M_1$，则$M=De_1\oplus M_1$，$Dg\cap N_1=\{0\}$。
任取$x=\sum\limits_{i=1}^{n}b_if_i\in N$，则$b_i\in I_i$，故$b_i=p^{k_i}b_i'$，于是$x-b_1g\in N\cap M_1=N_1$，即
\begin{align*}
N=Dg\oplus N_1,\quad g=p^{k_1}e_1.
\end{align*}
又$r(M_1)=n-1$，且
\begin{align*}
M_1/N_1\cong (N+M_1)/N\subseteq M/N,
\end{align*}
故$M_1/N_1$为$p$-模。由归纳假设，存在$M_1$的基$e_2,e_3,\cdots,e_n$，使得$p^{k_2}e_2,\cdots,p^{k_n}e_n$为$N_1$的基。
因此存在$M$的基$e_1,e_2,\cdots,e_n$，使得$p^{k_1}e_1,\cdots,p^{k_n}e_n$为$N$的基。

\end{proof}

\begin{theorem}
设$M'$是$D$上有限生成的$p$-模，则
\begin{enumerate}[(1)]
\item $M'$可分解为若干循环$p$-模的直和，即$\exists M_1,M_2,\cdots,M_m$使得$M_i=Dy_i$，满足
\begin{align*}
M'=\bigoplus\limits_{i=1}^{m}M_i,\quad \mathrm{ann}y_i=\langle p^{k_i}\rangle\ (k_i\geqslant1);
\end{align*}
\item 若分解满足$k_1\leqslant k_2\leqslant\cdots\leqslant k_m$，则数组$k_1,k_2,\cdots,k_m$由$M'$唯一确定。
\end{enumerate}
\end{theorem}
\begin{proof}

(1) 设$x_1,x_2,\cdots,x_n$为$M'$的生成元，$M$为$D$上秩为$n$的自由$D$-模。
由自由模的同态性质，存在$M$到$M'$的满同态$\eta$，满足$\eta(f_i)=x_i$，其中$f_1,\cdots,f_n$为$M$的基。
令$N=\ker\eta$，则$M/N\cong M'$为$p$-模。由引理5.4.3，存在$M$的基$e_1,\cdots,e_n$及数组$0\leqslant m_1\leqslant m_2\leqslant\cdots\leqslant m_n$，使得$N$的基为$p^{m_1}e_1,\cdots,p^{m_n}e_n$。
于是
\begin{align*}
M/N\cong\bigoplus\limits_{i=1}^{n}De_i/Dp^{m_i}e_i.
\end{align*}
若$m_i=0$，则$De_i/Dp^{m_i}e_i=\{0\}$。将$m_1,\cdots,m_n$中非零项记为$k_1,\cdots,k_m$，满足$1\leqslant k_1\leqslant k_2\leqslant\cdots\leqslant k_m$，则
\begin{align*}
\bigoplus\limits_{i=1}^{m}De_{i+n-m}/Dp^{k_i}e_{i+n-m}\cong M/N\cong M'.
\end{align*}
注意到$\mathrm{ann}(e_{i+n-m}+Dp^{k_i}e_{i+n-m})=\langle p^{k_i}\rangle$，故$M'$可分解为$M'=\bigoplus\limits_{i=1}^{m}Dy_i$，其中$\mathrm{ann}y_i=\langle p^{k_i}\rangle$。

(2) 对$m$用数学归纳法证明唯一性。
当$m=0$时显然成立。设$m-1$时成立，下证$m$时成立。
设$M'$另有分解
\begin{align*}
M'=\bigoplus\limits_{j=1}^{r}Dz_j,\quad \mathrm{ann}z_j=\langle p^{l_j}\rangle,\quad l_1\leqslant l_2\leqslant\cdots\leqslant l_r.
\end{align*}
由
\begin{align*}
z_r=\sum\limits_{i=1}^{m}a_iy_i,\quad y_m=\sum\limits_{j=1}^{r}b_jz_j,
\end{align*}
得$p^{k_m}z_r=p^{l_r}y_m=0$，故$k_m=l_r$。
若$\exists j$满足$\mathrm{ann}z_j=\langle p^{k_m}\rangle$且$p\mid b_j$，则$p^{k_m-1}y_m=0$，与$\mathrm{ann}y_m=\langle p^{k_m}\rangle$矛盾。
故$\exists j\leqslant r$使得$l_j=l_r=k_m$且$p\nmid b_j$。不妨设$j=r$，即$(p^{k_m},b_r)=1$。
于是$\exists u,v\in D$使得$up^{k_m}+vb_r=1$，故
\begin{align*}
vy_m&=\sum\limits_{j=1}^{r-1}vb_jz_j+vb_rz_r\\
&=\sum\limits_{j=1}^{r-1}vb_jz_j+(1-up^{k_m})z_r\\
&=z_r+\sum\limits_{j=1}^{r-1}vb_jz_j,
\end{align*}
即
\begin{align*}
z_r=\sum\limits_{j=1}^{r-1}(-vb_j)z_j+vy_m.
\end{align*}
任取$x\in\left(\bigoplus\limits_{j=1}^{r-1}Dz_j\right)\cap Dy_m$，则
\begin{align*}
x=\sum\limits_{j=1}^{r-1}c_jz_j=cy_m=\sum\limits_{i=1}^{m}cb_iz_i,
\end{align*}
比较系数得$cb_rz_r=0$。由$\mathrm{ann}z_r=\langle p^{l_r}\rangle$得$p^{l_r}\mid cb_r$，结合$p\nmid b_r$得$p^{l_r}\mid c$。
又$\mathrm{ann}y_m=\langle p^{l_r}\rangle$，故$x=0$，因此
\begin{align*}
M'=\bigoplus\limits_{j=1}^{r-1}Dz_j\oplus Dy_m.
\end{align*}
于是
\begin{align*}
M'/Dy_m\cong\bigoplus\limits_{i=1}^{m-1}Dy_i,\quad M'/Dy_m\cong\bigoplus\limits_{j=1}^{r-1}Dz_j.
\end{align*}
由归纳假设，$m-1=r-1$且$l_j=k_j,\ 1\leqslant j\leqslant m-1$，唯一性得证。

\end{proof}

\begin{corollary}
设$N$为秩$n$自由$D$-模$M$的子模，$M/N$为$p$-模。若$e_1,\cdots,e_n$与$f_1,\cdots,f_n$均为$M$的基，$p^{k_1}e_1,\cdots,p^{k_n}e_n$与$p^{l_1}f_1,\cdots,p^{l_n}f_n$均为$N$的基，且$k_1\leqslant\cdots\leqslant k_n$，$l_1\leqslant\cdots\leqslant l_n$，则$l_i=k_i,\ \forall i$。
\end{corollary}
\begin{proof}

由
\begin{align*}
M/N\cong\bigoplus\limits_{i=1}^{n}De_i/Dp^{k_i}e_i\cong\bigoplus\limits_{i=1}^{n}Df_i/Dp^{l_i}f_i,
\end{align*}
结合定理5.4.3的(2)得$l_i=k_i$。

\end{proof}

\begin{theorem}
设$D$为主理想整环，$M$为$D$上有限生成模，则$M$可分解为有限个循环子模的直和，即
\begin{align*}
M=Dx_1\oplus Dx_2\oplus\cdots\oplus Dx_n,
\end{align*}
且$\mathrm{ann}x_i$为$\langle0\rangle$或$\langle p_i^{k_i}\rangle$（$1\leqslant i\leqslant n$，$p_i$为$D$的素元）。
记$M$的秩为自由子模的个数$|\{x_i\mid \mathrm{ann}x_i=\langle0\rangle\}|$，$\{p_i^{k_i}\}$为$M$的\textbf{初等因子}，则$M$的秩与初等因子由$M$唯一确定。
\end{theorem}
\begin{proof}

结合主理想整环上有限生成模的自由-扭模分解、扭模的素分支分解、$p$-模的循环模直和分解，直接可得上述结论，唯一性由各步分解的唯一性保证。

\end{proof}

\begin{lemma}
设$D$为主理想整环，$M$为$D$-模，$x,y\in M$，满足$\mathrm{ann}x=\langle a\rangle$，$\mathrm{ann}y=\langle b\rangle$，且$(a,b)=1$，则
\begin{align*}
D(x+y)=Dx\oplus Dy,\quad \mathrm{ann}(x+y)=\langle ab\rangle.
\end{align*}
\end{lemma}
\begin{proof}

由$\mathrm{ann}x=\langle a\rangle,\mathrm{ann}y=\langle b\rangle,(a,b)=1$，结合引理5.4.1的(4)得
\begin{align*}
Dx\cap Dy\subseteq M(a)\cap M(b)=M((a,b))=M(1)=\{0\}.
\end{align*}
又$D(x+y)\subseteq Dx\oplus Dy$。取$u,v\in D$使得$ua+vb=1$，则
\begin{align*}
x=(1-ua)x=vbx=vb(x+y),\quad y=(1-vb)y=uay=ua(x+y),
\end{align*}
故$Dx\oplus Dy\subseteq D(x+y)$，即$D(x+y)=Dx\oplus Dy$。

又$ab(x+y)=bax+aby=0$，故$ab\in\mathrm{ann}(x+y)$。
若$c(x+y)=0$，则$cx\in Dx,cy\in Dy$，由$cx+cy=0$得$cx=cy=0$，故$a\mid c,b\mid c$。
因$(a,b)=1$，故$ab\mid c$，即$\mathrm{ann}(x+y)=\langle ab\rangle$。

\end{proof}

\begin{theorem}
设$D$为主理想整环，$M$为有限生成$D$-模，则$M$可分解为循环子模的直和
\begin{align*}
M=\bigoplus\limits_{i=1}^{s}Dz_i,
\end{align*}
满足
\begin{align*}
\mathrm{ann}z_1\supseteq\mathrm{ann}z_2\supseteq\cdots\supseteq\mathrm{ann}z_s,\quad \mathrm{ann}z_1\neq D,
\end{align*}
该分解由$M$唯一确定。记$\mathrm{ann}z_i=\langle d_i\rangle$，称$d_i$为$M$的\textbf{不变因子}。
\end{theorem}
\begin{proof}

由定理5.4.4，$M$有初等因子分解
\begin{align*}
M=\bigoplus\limits_{i=0}^{s}\bigoplus\limits_{j=1}^{\lambda_i}Dx_{ij},
\end{align*}
其中$\mathrm{ann}x_{0j}=\langle0\rangle$，$\mathrm{ann}x_{ij}=\langle p_i^{k_{ij}}\rangle$，$p_i$为互不相伴素元，且$k_{i1}\geqslant k_{i2}\geqslant\cdots\geqslant k_{i\lambda_i}$。
约定$j>\lambda_i$时$x_{ij}=0$，令
\begin{align*}
x^{(j)}=\sum\limits_{i=1}^{s}x_{ij},
\end{align*}
由引理5.4.4得
\begin{align*}
Dx^{(j)}=\bigoplus\limits_{i=1}^{s}Dx_{ij},\quad \mathrm{ann}x^{(j)}=\langle p_1^{k_{1j}}p_2^{k_{2j}}\cdots p_s^{k_{sj}}\rangle.
\end{align*}
因此
\begin{align*}
M=\left(\bigoplus\limits_{j}Dx^{(j)}\right)\oplus\left(\bigoplus\limits_{j=1}^{\lambda_0}Dx_{0j}\right),
\end{align*}
且
\begin{align*}
\langle0\rangle=\mathrm{ann}x_{0j}\subseteq\mathrm{ann}x^{(1)}\subseteq\mathrm{ann}x^{(2)}\subseteq\cdots\subseteq\mathrm{ann}x^{(\lambda)},
\end{align*}
其中$\lambda=\max\{\lambda_1,\lambda_2,\cdots,\lambda_s\}$。

令$s=\lambda+\lambda_0$，将$x_{01},\cdots,x_{0\lambda_0},x^{(1)},\cdots,x^{(\lambda)}$记为$z_s,z_{s-1},\cdots,z_1$，则
\begin{align*}
M=Dz_1\oplus Dz_2\oplus\cdots\oplus Dz_s,
\end{align*}
且
\begin{align*}
\mathrm{ann}z_1\supseteq\mathrm{ann}z_2\supseteq\cdots\supseteq\mathrm{ann}z_s.
\end{align*}
由$\lambda_0,p_i,k_{ij}$的唯一性，$\mathrm{ann}z_i=\langle d_i\rangle$由$M$唯一确定。

\end{proof}



\end{document}